\section{Evaluation Data and Methodological Evidence}
\label{sec:evaluation}

This section synthesizes the literature that determines how evidence about
accent and speech translation can be interpreted. It first examines dataset
and reference design, then considers translation metrics and statistical
reliability. Together, these themes establish the controls required for a
credible comparison without turning the survey into an experimental procedure.

\subsection{Accent-Labelled Speech Data}

Common Voice provides a practical source of accent-labelled speech, but the
status of those labels needs to be stated carefully. Ardila et al. describe
\texttt{accent} as one of three optional demographic fields self-reported by
speakers, alongside age and gender; the paper does not report expert phonetic
verification of this field \cite{ardila2020commonvoice}. Prabhu et al. later
extract an English validated subset, \texttt{MCV\_ACCENT}, and organize 14
Common Voice accent categories into five seen accents used for training and
validation and nine unseen accents used for testing. Their train, development,
and test sets are speaker-disjoint \cite{prabhu2023accentcodebooks}. This work
demonstrates that Common Voice metadata can support controlled
accent-robustness experiments, but neither paper establishes independent
validation of the accent categories. A version caveat is also necessary:
Ardila et al. describe Common Voice data from 2019, whereas Prabhu et al. do
not report the exact Common Voice release used for \texttt{MCV\_ACCENT}; their
counts and split details therefore should not be treated as describing an
identical dataset snapshot.

\subsection{Speech-Translation Data and Reference Quality}

CoVoST 2 supplies the translation layer that Common Voice itself lacks. Wang
et al. derive it from validated Common Voice read speech and send deduplicated
source transcripts, without the corresponding audio, to professional
translators \cite{wang2021covost2}. Translation quality is checked using
language-model perplexity, LASER scores, and a length-ratio heuristic.
Low-scoring items are manually inspected and retranslated, while sentences
that cannot be translated properly, often because of missing context, are
marked in the release. This process provides documented professional
translation and quality-control provenance. The paper also reports cascaded
and end-to-end ST baselines on the corpus; its end-to-end ST category
corresponds to the Direct ST category used in this survey
\cite{wang2021covost2}. However, the references are conditioned on text rather
than audio, and the authors warn that the CoVoST multi-speaker development and
test splits contain repeated sentences that may skew BLEU or WER. These
properties must therefore be considered when designing a comparable
evaluation.

\subsection{Dataset Design for Controlled Accent Comparisons}

Together, the three papers show that a fair comparison between Direct and
Cascaded ST requires careful control of the evaluation data. Common Voice
keeps speakers separate across its official training, development, and test
sets and removes repeated text sentences from those splits
\cite{ardila2020commonvoice}. CoVoST 2 adds professional translation
references and reports both cascaded and end-to-end ST baselines, but it was
not designed to isolate accent effects, and its alternative multi-speaker
development and test splits contain repeated sentences
\cite{wang2021covost2}. Prabhu et al. provide a stronger accent-control
template: speaker-disjoint seen/unseen accent splits and an explicit
transcript-overlap policy. Their curation deliberately permits some transcript
overlap across splits, while keeping most development and test examples
disjoint from training in both speaker and transcript
\cite{prabhu2023accentcodebooks}.

Consequently, a pipeline difference should be evaluated on matched utterances
and references, with speaker identity, transcript overlap, accent imbalance,
dataset version, and recording-condition variation controlled or at
least reported. These controls improve comparability but do not establish
architecture as the sole cause: the evaluated systems may also differ in
training data, capacity, objectives, or decoding. Once the data are comparable,
the metric and statistical-reliability literature determines how strongly any
observed difference can be interpreted.

\subsection{Evaluation Metrics}

WER only signifies transcription error when compared to a reference sentence.
It helps as a diagnostic metric (similar to WIL) in the Cascaded approach,
where a gap in the ASR stage can be identified. Moreover, this is not available
in Direct ST, as there is no intermediary layer available for that approach.
As a result, it helps to identify faults in one approach but has no impact on
the evaluation of the other. Nor does WER show how much translation quality is
lost: recognition errors of similar size can produce very different amounts of
translation damage \cite{le2017disentangling}. To adapt WER to multiple
references and differences in word selection and sequence, BLEU was derived on
top of it. BLEU attempts to address translation quality by taking additional
factors into consideration \cite{papineni2002bleu}.

BLEU builds on unigram precision by introducing modified precision, which caps
the maximum credit an output word can receive based on the reference sentence.
This penalises excessively repeated word outputs. It further extends
reliability by introducing a brevity penalty to penalise shorter outputs based
on the best length match. It achieved a significantly high correlation with
human judgement \cite{papineni2002bleu}. However, over a set of reference
outputs, a seemingly correct translation can score lower because BLEU does not
have access to the meanings of words. Moreover, it was tested over a
500-sentence corpus rather than using a sentence-level approach. The similarity
with human judgement holds at corpus level rather than for individual sentences
\cite{papineni2002bleu}. chrF is a character n-gram score which acts as a
singular tool for judging the overall quality of a translation because of its
tokenisation-independent and human-validation-based approach. The approach
puts $\beta$ times more emphasis on recall than on precision. When compared
across two $\beta$ values, $\beta = 3$ was found to have higher correlation
with human output than $\beta = 1$, which assigns equal weight to precision and
recall \cite{popovic2015chrf}. chrF++ is another metric which re-introduces
tokenisation into the mix by merging character n-grams with word unigrams and
bigrams. It was also validated against a direct-assessment approach, and
$\beta = 2$ was confirmed to perform best when assessed on English and Russian
targets \cite{popovic2017chrfpp}.

For comparison, BLEU scores are not comparable across various papers simply
because the parameters and variables that affect the value are not easily
accessible or clearly reported. As a result, direct comparison between BLEU
scores is similar to comparing measurements expressed in different units.
Differences can vary by up to 1.8 BLEU and average around 1.0
\cite{post2018clarity}. Differences in reference normalisation or tokenisation
cause variance that is difficult to ignore \cite{post2018clarity}. IWSLT 2021
shows a gap of 2 BLEU between the Cascaded and Direct systems with one reference
set and 0.7 with another \cite{anastasopoulos2021iwslt}. The 2 and 0.7 gaps are
of the same order of magnitude as Post's 1.0 to 1.8 variance
\cite{post2018clarity}. chrF was initially tested with only one non-European
language, Hindi, and its application to other writing systems, such as Chinese,
remains to be explored \cite{popovic2015chrf}. It remains unclear whether
chrF++ or chrF would be better because the additional benefits gained by
re-introducing word tokenisation for English-to-Chinese pairs have not been
derived. Popovi\'c submitted only plain chrF for Chinese because the concept of
a Chinese word is unclear \cite{popovic2017chrfpp}.

\subsection{Statistical Reliability and Controlled Evaluation}

A difference in BLEU score between two systems on one test set does not by
itself show that one system is better, so a statistically sound approach is
needed to reach this conclusion. However, because of the complex nature of the
metric itself, the bootstrap sampling method can be used to develop an
understanding of the scenario. Koehn validates the method by comparing systems
across many test sets of different sizes and checking how often its conclusions
hold \cite{koehn2004significance}. Bootstrap sampling is introduced, where a
BLEU score is given for a set of $n$ sentences. Repeatedly generating BLEU
scores about 1000 times for 300-sentence sets drawn with replacement from one
set of 300 sentences results in the remaining 95\% of BLEU scores falling in an
interval after the top and bottom 2.5\% are dropped
\cite{koehn2004significance}. When comparing two systems, paired bootstrap is
introduced to measure reliably whether one system outperforms the other. From
the same set of translated sentences, the evaluation metric is computed for
both systems 1000 times. A conclusion can be reached after identifying how
often one system wins over the other; if it wins over 95\% of the time, that
system can be said to be better with 95\% statistical significance
\cite{koehn2004significance}. For a Direct-versus-Cascaded comparison, scoring
both systems on the same test set is what makes this paired test applicable,
which matches the matched-utterance requirement of Section~V-C and the design
used by Bentivogli et al. \cite{bentivogli2021cascade}.

Koehn's experiment is based on the underlying assumption that each sentence
score is independent of the others. He states this for the analytical case, and
bootstrap resampling of sentences relies on it too
\cite{koehn2004significance}. Accent-labelled read speech breaks this
assumption. One speaker can contribute many clips, and the same source sentence
can appear several times \cite{ardila2020commonvoice,wang2021covost2}.
Resampling clips then treats linked observations as if they were separate, and
the uncertainty becomes smaller than it should be. Koehn also shows that the
factors that make sentences hard to translate occur in clusters, and clusters
widen the spread of scores: on 300 sentences taken in sequence, BLEU ranged
from 21 to 37 per cent, but on 300 sentences spread across the corpus it ranged
only from 27 to 31 per cent \cite{koehn2004significance}. A head-to-head
comparison of Direct and Cascaded ST on the same clips is the cleanest approach,
as the metric is computed at corpus level rather than at subset level, as is the
case with accent groups \cite{papineni2002bleu}. This is further supported by
Koehn's Table 3, which shows that small differences cannot be distinguished
even with 600 sentences \cite{koehn2004significance}. Any accent-level
difference between Direct and Cascaded ST can therefore be relied upon only
when these conditions are met, and not solely on the numerical gap between the
approaches.
